\section{Integración y conclusiones}
\label{sec:integracion}

\begin{table}[H]
  \centering
  \caption{Integración de las técnicas: estructura hallada, protagonistas y relación con las demás. Fuente: elaboración propia con UNSW-NB15.}\label{tab:integracion}
  \begin{tablacompacta}
  \footnotesize
  \setlength{\tabcolsep}{3pt}
  \begin{tabularx}{\linewidth}{@{}>{\raggedright\arraybackslash}p{1.45cm}>{\raggedright\arraybackslash}X>{\raggedright\arraybackslash}X>{\raggedright\arraybackslash}X@{}}
  \toprule
  Técnica & Estructura & Protagonistas & Coincidencias y discrepancias \\
  \midrule
  ACP (§\ref{sec:acp}) & 3 componentes, 74,5~\% de la varianza: CP1 magnitud del flujo frente a repetición (43,3~\%), CP2 volumen de tráfico, CP3 ritmo lento & CP1: \texttt{sjit}, \texttt{sinpkt} frente a las cinco \texttt{ct\_*}; CP2: \texttt{spkts}, \texttt{sbytes} & Aísla la repetición, como el AF y el cluster de variables; pero sus ejes son mixtos y no separan volumen de temporización \\
  AF (§\ref{sec:af}) & 3 factores (GLS, varimax), 68,6~\%: F1 Repetición, F2 Temporización, F3 Volumen & Cargas de 0,54 a 0,91, sin cargas cruzadas ambiguas & Coincide exactamente con el cluster de variables (ARI = 1,000) y con H1; ajuste global flojo (TLI 0,609, RMSEA 0,239) \\
  Cluster de variables (§\ref{sec:cluster}) & 3 bloques (Ward): temporización, volumen, repetición & Las cinco \texttt{ct\_*} juntas; \texttt{sinpkt}--\texttt{sjit} (0,85) & Igual al AF; el ACP mezcla los bloques en cada eje \\
  Cluster de observaciones (§\ref{sec:cluster}) & 6 grupos (k-means), silhouette 0,333, estructura débil & C1 transferencias voluminosas (1,4~\% ataque); C3 ráfagas repetidas sin respuesta (100~\%); C2 y C6 mal separados & Los grupos se ordenan en CP1 y F1 (C3: $-7{,}00$ y $+3{,}78$); en el ACM, C3 queda aislado en Dim 1. Con puntajes del ACP, ARI = 0,708 \\
  ACS y ACM (§\ref{sec:acm}) & ACS: tipo $\times$ estado (V = 0,405, 100~\% de inercia en 2 dim.). ACM: 2 dim., 41,4~\% crudo (87,4~\% ajustado) & Dim 1: INT, sin respuesta, duración corta y UDP frente a FIN y TCP; Dim 2: DNS/CON frente a protocolo \emph{otros} & Dim 1 reproduce el contraste de CP1 y el cluster 3; \texttt{tipo\_trafico} casi no pesa ($\eta^2$ = 0,242) \\
  Logística (§\ref{sec:logistica}) & M2 con PC1--PC3 y categorías: AUC 0,882 en prueba (IC 95~\% 0,864--0,900) & PC1 (OR 5,01), PC3 (2,54), PC2 (0,59); INT, CON\_otros (0,19), dns & Coincide con el ACM en INT y CON\_otros; discrepa en dns (signo opuesto al bruto) y los OR categóricos son inestables \\
  \bottomrule
  \end{tabularx}
  \end{tablacompacta}
\end{table}

\textbf{Hipótesis.}
\begin{singlespace}
\begin{itemize}\setlength{\itemsep}{0pt}\setlength{\topsep}{1pt}\setlength{\parsep}{0pt}
  \item \textbf{H1, confirmada.} El AF y el cluster de variables dan los mismos tres bloques (ARI = 1,000; cargas $\ge 0{,}54$). El ACP pide también tres componentes, pero no los recupera como ejes puros (secciones~\ref{sec:acp} y~\ref{sec:af}).
  \item \textbf{H2, parcialmente confirmada.} El cluster 1 es 98,6~\% normal y el 3 es 100~\% ataque, pero los otros cuatro tienen entre 34,9 y 72,6~\% de ataques. Además, el cluster 3 sí distingue un tipo (130 de 147 flujos son Generic): falla la parte de ``no distinguen tipos'' (sección~\ref{sec:cluster}).
  \item \textbf{H3, parcialmente confirmada.} INT y sin respuesta se asocian con Otros ataques (62,3~\% y 63,2~\% frente a 5,9~\% y 6,9~\% de Normal), pero Exploits es casi todo FIN con respuesta, y dns y http no son exclusivos de lo normal (sección~\ref{sec:acm}).
  \item \textbf{H4, confirmada.} AUC de prueba 0,882, con límite inferior 0,864 $>$ 0,80; solo con las componentes (M1) sería 0,795 (sección~\ref{sec:logistica}).
\end{itemize}
\end{singlespace}

\textbf{Bloques y clusters.} Los bloques coinciden entre AF y cluster de variables, y el ACP los recupera solo en parte: separa la repetición, pero mezcla volumen y temporización en CP1. Los clusters se ordenan sobre todo por CP1 ($\eta^2 = 0{,}86$; CP2 y CP3, 0,65) y por los factores (F2 0,78, F3 0,73, F1 0,72): por ejemplo, el cluster 3 tiene F1 = $+3{,}78$ (repetición) y el 2, F2 = $+0{,}88$ (sesiones largas). Es en parte circular, porque se agrupó con las mismas variables. Las categorías del ACM lo confirman con variables que no entraron al cluster: el cluster 3 es 100~\% INT y sin respuesta (Dim 1 = 2,16) y el 5 es breve y UDP (Dim 1 = 1,37), mientras que los clusters 1, 2 y 4 se superponen en TCP/FIN y solo se distinguen por duración y bytes.

\textbf{Sorpresas.} (i) La repetición de conexiones, que esperábamos asociada a los ataques, no los distingue por sí sola: los clusters 3 y 6 la tienen alta, pero uno es 100~\% ataque y el otro 34,9~\%. (ii) El cluster 3 aísla ataques Generic. (iii) INT es redundante con la falta de respuesta (913 de 913 flujos), así que el ACM y la logística no pueden separarlos. (iv) El caso Heywood: \texttt{log\_spkts} supera comunalidad 1 con casi cualquier estimador y solo GLS lo evita. (v) Los OR categóricos son inestables entre particiones (INT de 93 a 375, udp de 0,42 a 0,14). (vi) El ACP, con ejes mixtos, resume más varianza (74,5~\%) pero el AF, con bloques puros, interpreta mejor.

\textbf{Limitaciones.} Los ataques son sintéticos, de laboratorio; sin duplicados, las proporciones no son prevalencias y la probabilidad logística es relativa a esta base (52~\% de ataques); excluimos TTL por artefacto; los \texttt{ct\_*} usan ventanas de 100 conexiones, así que los flujos no son independientes; con $n = 5.000$ Bartlett rechaza cualquier asociación y pesa más el KMO; persiste asimetría pese al logaritmo; el ajuste global del AF es flojo; k = 6 tiene una silhouette casi empatada con k = 2 a 5; los terciles del ACM son un criterio estadístico y no sustantivo; hay colinealidad entre categóricas en la logística (VIF $\approx$ 17--18); usamos una sola partición 70/30; y los datos son observacionales, por lo que describimos asociaciones y no causas.

\textbf{Próximos pasos.} Validar con el testing set oficial procesado igual (mismo logaritmo, cortes y ACP), repetir con validación cruzada, y probar modelos con interacciones o regularización que reduzcan la inestabilidad de los OR.

% REVISAR GRUPO: opinión propia
\textbf{Conclusión.} Respondemos a la pregunta guía así: los flujos se organizan en tres patrones, volumen, temporización y repetición, y esos patrones, junto con el estado, el servicio y el protocolo, permiten estimar la probabilidad de ataque con un AUC de 0,882 en prueba, aunque con 28,7~\% de los normales marcados como ataque. A nuestro juicio (valoración del grupo), lo más sólido es la estructura de variables, que tres técnicas recuperan, y lo más frágil es la lectura de grupos y de odds ratios aislados. Creemos que el estado de la conexión aporta más que cualquier patrón numérico, pero no podemos extrapolar a una red real.
